\label{p3-20260826:chap:medida-accion-positiva}
\index{mesure!catalogue d’émissions}
\index{action!cocycle positif}
\index{trace normalisée}

La section précédente identifie isométriquement un secteur positif de rang trois dans
l’interface orientée \(t=7\), la dodécaphase et le module d’incidence de Witt. La construction
suivante dérive l’autre coefficient de la loi angulaire, \(5/468\), et
démontre pourquoi les deux termes se combinent par une somme positive. La
question exige de séparer
trois opérations :
\begin{enumerate}
  \item compter les présentations microscopiques ;
  \item compter les émissions distinctes après formation du quotient TPK ;
  \item attribuer une action additive à une route.
\end{enumerate}
Les deux premières opérations produisent des mesures différentes. La troisième se
caractérise par une propriété universelle et non par la proximité ultérieure
d’un nombre décimal.

Nous noterons \(\Gtr\) le graphe orienté à \(1944\) états du lecteur de
chemins enrichis. Sa catégorie libre de chemins sera notée
\(\mathsf P(\Gtr)\). La définition locale de ses arêtes et de son action est
développée dans la section correspondante.

\section{Espace de germes, quotient d’émissions et fibres régionales}

Soit
\[
 \Omega_{\rm seed}
 =\bigl((\mathbb Z/9\mathbb Z)^2\times D_4\bigr)^2,
 \qquad
 |\Omega_{\rm seed}|=104\,976,
\]
l’ensemble exécutable des paires de germes de l’émetteur fini, et soit
\[
 F:\Omega_{\rm seed}\twoheadrightarrow\mathcal U_6,
 \qquad
 |\mathcal U_6|=468,
\]
l’application qui conserve l’émission sextuple distincte. Le recensement exact des
fibres est
\[
\begin{array}{c|c|c}
|F^{-1}(U)|&\#\{U\text{ de cette taille}\}&\text{germes apportés}\\
\hline
144&432&62\,208\\
432&18&7\,776\\
1\,944&18&34\,992.
\end{array}
\]
Les trois contributions ont pour somme \(104\,976\). Par conséquent, \(F\) n’est pas un
revêtement régulier : des émissions différentes possèdent des nombres différents de
présentations microscopiques.

La microfibre de \(\pi\) contient cinq émissions distinctes. Quatre ont
pour multiplicité \(144\) et une a pour multiplicité \(432\). Cette donnée suffit
à démontrer que la mesure uniforme sur les germes et la mesure uniforme
sur les émissions ne sont pas identiques.

\section{\(2\) mesures exactes sur le quotient des émissions et des mots}

Si chaque germe reçoit le poids \(1/104\,976\), la mesure image par \(F\) est
\[
 \nu_{\rm seed}(U)
 =\frac{|F^{-1}(U)|}{104\,976}.
\]
Ses atomes possibles sont
\[
 \frac1{729},\qquad\frac1{243},\qquad\frac1{54}.
\]
\Needspace{5\baselineskip}
En particulier,
\[
 \nu_{\rm seed}(\mathscr R_\pi^{\rm micro})
 =\frac{4\cdot144+432}{104\,976}
 =\boxed{\frac7{729}}.
\]

Si l’état élémentaire est, en revanche, une émission distincte
\(U\in\mathcal U_6\), l’algèbre diagonale des observables est
\[
 \mathcal D_{468}=\ell^\infty(\mathcal U_6)
\]
et son état de comptage normalisé est
\[
 \tau_{468}(f)
 =\frac1{468}\sum_{U\in\mathcal U_6}f(U).
\]
Soit \(\Pi_\pi^{(468)}\) la fonction caractéristique des cinq régions de la
microfibre. Alors
\[
 \boxed{\tau_{468}(\Pi_\pi^{(468)})=\frac5{468}.}
\]
La différence exacte entre les deux lectures est
\[
 \frac5{468}-\frac7{729}
 =\frac{41}{37\,908}.
\]
Ce ne sont pas deux approximations. La première mesure compte les présentations ; la
seconde compte les classes d’émission.

\begin{theorem}[Unicité de la mesure uniforme du catalogue]
\label{p3-20260826:thm:unicidad-medida-468}
L’unique probabilité sur \(\mathcal U_6\) invariante sous tout
réétiquetage de ses \(468\) atomes est la mesure uniforme.
De manière équivalente, l’unique trace positive et unitaire sur
\(M_{468}(\mathbb C)\) est
\[
 \tau_{468}(A)=\frac1{468}\operatorname{Tr}A.
\]
\end{theorem}

\begin{proof}
Les transpositions envoient tout atome sur tout autre. L’invariance
impose qu’ils aient tous la même masse, et la normalisation la fixe à
\(1/468\). Dans la version matricielle, l’invariance unitaire égalise la valeur
de tous les projecteurs de rang un. La décomposition spectrale et la
linéarité reconstruisent la trace normalisée.
\end{proof}

L’hypothèse décisive est manifeste : le lecteur travaille avec le
\emph{catalogue quotient}. Le théorème n’affirme pas que la mesure uniforme des
germes devienne uniforme lorsqu’on oublie les fibres. La symétrie utilisée
est la naturalité de l’objet fini sans étiquettes, et non une propriété
ergodique attribuée sans preuve à la dynamique microscopique.

\section{Hilbertisation canonique du quotient}

La relation entre germes et émissions peut être formulée sans ambiguïté.
Soient
\[
 H_{\rm seed}=\ell^2(\Omega_{\rm seed}),
 \qquad
 H_U=\ell^2(\mathcal U_6).
\]
L’émetteur détermine
\[
 \boxed{
 J_F\delta_U
 =|F^{-1}(U)|^{-1/2}
 \sum_{\omega\in F^{-1}(U)}\delta_\omega.
 }
\]
Les fibres distinctes sont disjointes et chaque colonne a pour norme un ; par
conséquent,
\[
 J_F^*J_F=I_{H_U}.
\]
De plus, pour tout observable \(f\) du catalogue,
\[
 M_{f\circ F}J_F=J_FM_f.
\]
L’image \(J_FH_U\) est exactement le sous-espace des fonctions constantes
sur chaque fibre.

Si \(F^\sharp:H_U\to H_{\rm seed}\) désigne l’application induite par
précomposition, sans normalisation,
alors
\[
 J_F
 =F^\sharp\bigl((F^\sharp)^*F^\sharp\bigr)^{-1/2}.
\]
C’est le facteur isométrique de la décomposition polaire de cette application. Cette
expression démontre sa naturalité sous les isomorphismes et les réétiquetages de
l’émetteur. Elle n’implique pas une fonctorialité fausse sous toute composition de
surjections non régulières : la normalisation dépend des fibres de
l’application concrète.

Dans le langage des mesures, le relèvement de la trace du quotient qui est
uniforme à l’intérieur de chaque fibre vaut
\[
 \widetilde\tau(\omega)
 =\frac1{468\,|F^{-1}(F\omega)|}.
\]
Chaque fibre reçoit une masse totale \(1/468\). La compensation inverse par
la multiplicité, et non une invocation informelle des orbites, constitue la preuve exacte.

\begin{corollary}[Canonicité de \(5/468\)]
\label{p3-20260826:cor:canonicidad-cinco-468}
Une fois le lecteur régional défini sur les émissions distinctes, la masse
de la microfibre de \(\pi\) est imposée :
\[
 \tau_{468}(\Pi_\pi^{(468)})
 =\frac{\operatorname{rank}\Pi_\pi^{(468)}}{468}
 =\frac5{468}.
\]
\end{corollary}

\section{Défaut positif entre les mesures régionale et dodécaphasique}

Écrivons
\[
 H_{\rm ang}=H_U\otimes V_{11}.
\]
Le scalaire interne
\[
 \Delta_4
 =\frac{\pi-e\log\pi}{270}
 \left(1+\frac\pi{729}\right)
 =0.0001111964226921402\ldots
\]
est construit avec les lectures HMT déjà fixées de \(\pi\) et \(e\), et avec les
entiers \(270\) et \(729\) du registre. Il constitue ici une entrée mathématique
prédimensionnelle : ni unités ni observables physiques n’interviennent.

La construction interne du scalaire ne détermine pas à elle seule quel opérateur
il doit multiplier. Pour garder cette prémisse visible, nous déclarons :
\begin{description}
  \item[\HDef --- défaut de changement de feuillet.] La coordonnée dodécaphasique du défaut
  de changement de feuillet est \(\Delta_4P_3\).
\end{description}
\HAdd fixera la forme additive de l’action, mais ne sélectionne pas \HDef : pour tout
\(c\ge0\), le même argument admettrait la paire
\((\Pi_\pi^{(468)},cP_3)\). Démontrer \HDef directement à partir d’une transition
totale du noyau topologique de phase enrichi est une condition supplémentaire.

Nous définissons la somme de Kronecker
\[
 \boxed{
 D_{\rm switch}^{\rm dod}
 =\Pi_\pi^{(468)}\otimes I_{11}
 +\Delta_4 I_{468}\otimes P_3.
 }
\]
Les deux termes sont positifs, commutent et agissent sur des facteurs différents.
Par conséquent,
\[
 D_{\rm switch}^{\rm dod}\ge0.
\]
L’entrelaceur \(U_W\) du
\cref{thm:entrelazador-dodecafase-witt} produit la carte de Witt :
\[
\begin{aligned}
 D_{\rm switch}^{\rm Witt}
 &=(I_{468}\otimes U_W)
 D_{\rm switch}^{\rm dod}
 (I_{468}\otimes U_W)^*\\
 &=\Pi_\pi^{(468)}\otimes I_{E_{30}}
 +\Delta_4I_{468}\otimes Q_3.
\end{aligned}
\]
Ainsi, le même opérateur positif se lit dans la dodécaphase et dans le module
d’incidence.

La chaîne isométrique orientée identifie également les secteurs actifs
\[
 I_{L_{\rm or}}\otimes P_G,\qquad P_3,\qquad Q_3.
\]
On n’affirme pas une équivalence des compléments complets des trois
modules : \(W_{GK}\) est une équivalence partielle des rangs
sélectionnés. C’est exactement l’information nécessaire pour identifier
la seconde contribution de la trace.

\begin{theorem}[Portée exacte de l’entrelacement et de l’induction de l’action]
\label{p3-20260826:thm:alcance-induccion-accion}
Soient \(W_{GK}\) et \(U_W\) les isométries partielles du
\cref{thm:identificacion-isometrica-tres-proyectores}. Alors :
\begin{enumerate}
  \item pour tout \(c\geq0\),
  \[
   W_{GK}^{*}(cP_3)W_{GK}
   =c(I_{L_{\rm or}}\otimes P_G),
   \qquad
   U_W(cP_3)U_W^{*}=cQ_3;
  \]
  par conséquent, l’entrelacement transporte le coefficient, mais ne le
  sélectionne pas ;
  \item si \(P_{\rm sgn}\) et \(P_{\rm std}\) sont les projecteurs centraux
  de la décomposition
  \[
   P_3V_{11}\big|_{H_*}
   \simeq \operatorname{sgn}_{H_*}\oplus\operatorname{std},
   \qquad
   \operatorname{rank}P_{\rm sgn}=1,
   \quad
   \operatorname{rank}P_{\rm std}=2,
  \]
  tout opérateur positif autoadjoint supporté sur \(P_3V_{11}\) et
  \(H_*\)-équivariant a la forme
  \[
   D_{a,b}=aP_{\rm sgn}+bP_{\rm std},
   \qquad a,b\geq0.
  \]
  Même l’imposition de la trace de \(\Delta_4P_3\) laisse la famille
  \[
   a+2b=3\Delta_4.
  \]
  L’isotropie \(a=b\) ne se déduit pas de la symétrie résiduelle \(S_3\) ;
  \item la probabilité uniforme du catalogue,
  \(\mu_U(U)=1/468\), est la trace canonique du quotient fini et son
  relèvement compensé est celui que détermine \(J_F\). Ce n’est pas la
  mesure image de la probabilité uniforme sur les germes, dont les atomes sont
  \(1/729,1/243,1/54\). En particulier,
  \[
   \mu_U(\mathscr R_\pi^{\rm micro})=\frac5{468},
   \qquad
   F_\#\mu_{\rm seed}(\mathscr R_\pi^{\rm micro})=\frac7{729};
  \]
  \item la loi complète
  \[
   \delta_*=\frac5{468}+\frac3{11}\Delta_4,
   \qquad
   \lambda_*=1-\delta_*,
   \qquad
   a_*(e)=g(e)+\lambda_*s(e)
  \]
  est un théorème relatif aux cinq clauses \HTr--\HAdd--\HDef--\HRes--\HLoc. Le moteur
  exécutable du noyau topologique de phase réduit n’induit pas, par revêtement orienté, l’opérateur ramifié
  \Gtr : la dynamique associée revient à l’identité après \(54\) itérations,
  tandis que chaque état de
  \Gtr possède quatre successeurs distincts et aucune boucle.
\end{enumerate}
Par conséquent, les résultats précédents démontrent la canonicité des identifications
isométriques, de la trace sur le catalogue quotient et du caractère additif
dans le cadre de leurs prémisses. L’affirmation plus forte selon laquelle une transition totale
du noyau topologique de phase enrichi induit nécessairement \(\Delta_4P_3\), la mesure uniforme et la
loi locale exige encore le théorème dynamique supplémentaire spécifié dans
l’analyse de l’induction.
\end{theorem}

\begin{proof}
Les identités du premier point s’obtiennent en multipliant
\(W_{GK}^{*}W_{GK}=I_{L_{\rm or}}\otimes P_G\),
\(W_{GK}W_{GK}^{*}=P_3\) et \(U_WP_3U_W^{*}=Q_3\) par le scalaire \(c\).
Pour le second point, la restriction à \(H_*\simeq S_3\) est une somme sans
multiplicités de deux représentations irréductibles non isomorphes. Le lemme de
Schur diagonalise tout opérateur autoadjoint équivariant dans ces deux blocs ;
la positivité équivaut à \(a,b\geq0\). Sa trace normalisée est
\((a+2b)/11\), ce qui prouve la liberté résiduelle même après fixation de la
trace. L’action complète de \(A_5\) sur le triplet irréductible réduirait
la famille à \(cP_3\), mais ne fixerait pas non plus la valeur de \(c\).

Le troisième point découle du recensement des fibres : \(F\) a pour multiplicités
\(144,432,1944\), de sorte que la mesure image pondère une émission par son
nombre de présentations. La trace \(\tau_{468}\), en revanche, attribue le
même poids à chaque classe ; \(J_F\) réalise exactement son relèvement
normalisé. Le quatrième point est l’enchaînement logique déjà démontré :
\HTr fixe la trace du quotient, \HAdd la somme des traces, \HDef déclare
\(\Delta_4P_3\), \HRes définit la capacité résiduelle et \HLoc la porte sur les arêtes.
L’énumération finie prouve l’obstruction entre le noyau réduit et
\Gtr : un revêtement orienté préserverait l’étoile sortante et une
semi-conjugaison du retour identité exigerait des boucles dans l’image,
en contradiction avec les deux recensements de \Gtr.
\end{proof}

Le spectre complet de \(D_{\rm switch}^{\rm dod}\) est
\[
\begin{array}{c|c}
\text{valeur propre}&\text{multiplicité}\\ \hline
0&(468-5)(11-3)=3704\\
\Delta_4&(468-5)3=1389\\
1&5(11-3)=40\\
1+\Delta_4&5\cdot3=15.
\end{array}
\]
Par conséquent,
\[
\begin{aligned}
 \tau_{468\cdot11}(D_{\rm switch})
 &=\tau_{468}(\Pi_\pi^{(468)})+\Delta_4\tau_{11}(P_3)\\
 &=\boxed{\frac5{468}+\frac3{11}\Delta_4}.
\end{aligned}
\]
La même identité vaut dans la carte de Witt. Dans la carte du secteur orienté, le terme
\(3/11\) est la trace de \(P_G\) sur l’espace homogène à onze branches.

\section{Propriété universelle de la fonctionnelle positive d’action}

Le mot \emph{action} est employé au sens additif : la concaténation de deux
routes reçoit la somme de leurs incréments. Pour séparer ce contenu d’une
formule choisie, nous définissons d’abord la catégorie dans laquelle la loi sera unique.

\begin{definition}[Caractère positif d’action HMT]
\label{p3-20260826:def:caracter-accion-hmt}
Un caractère local d’action sur le produit
régional--dodécaphasique est une application
\[
 \delta:
 M_{468}(\mathbb C)_+\oplus M_{11}(\mathbb C)_+
 \longrightarrow\mathbb R_{\ge0}
\]
qui satisfait :
\begin{enumerate}
  \item additivité exacte dans le cône produit ;
  \item homogeneidad positiva;
  \item invariance par conjugaison unitaire indépendante dans les deux
  facteurs ;
  \item normalisation bifactorielle,
  \[
  \delta(I_{468},0)=\delta(0,I_{11})=1.
  \]
\end{enumerate}
\end{definition}

L’additivité traduit la concaténation ; l’invariance élimine les
coordonnées ; la normalisation bifactorielle fixe une échelle dans chaque terme.
Il ne s’agit pas d’une fonctionnelle unitaire sur l’algèbre somme directe : son unité
est \((I_{468},I_{11})\) et l’additivité donne
\(\delta(I_{468},I_{11})=2\). Aucune clause ne contient la valeur d’un
angle.

\begin{theorem}[Unicité du caractère positif et capacité résiduelle relative]
\label{p3-20260826:thm:ley-accion-hmt}
Dans la classe de la définition précédente, il existe un unique caractère positif
d’action HMT\@. C’est
\[
 \boxed{\delta(A,B)=\tau_{468}(A)+\tau_{11}(B).}
\]
Appliqué au défaut déclaré par \HDef,
\((\Pi_\pi^{(468)},\Delta_4P_3)\), il produit
\[
 \delta_*=\frac5{468}+\frac3{11}\Delta_4.
\]
Si l’on déclare en outre \HRes, selon laquelle la capacité autoduale de changement de
feuillet est normalisée à un et le défaut la réduit par
\(\lambda=1-\delta\), la capacité résiduelle dans le cadre de \HTr--\HAdd--\HDef--\HRes est
\[
 \boxed{
 \lambda_*
 =1-\delta_*
 =1-\frac5{468}-\frac3{11}\Delta_4.
 }
\]
\end{theorem}

\begin{proof}
L’additivité sépare les facteurs :
\[
 \delta(A,B)=\delta(A,0)+\delta(0,B).
\]
Chaque restriction est positive, homogène et invariante par conjugaison. Pour
un projecteur de rang un, l’invariance unitaire donne une valeur commune.
L’additivité sur une décomposition spectrale impose que chaque restriction
soit un multiple de la trace. La normalisation bifactorielle fixe ces multiples à
\(1/468\) et \(1/11\). Il en résulte la somme des traces. La dernière formule
applique la normalisation déclarée de la capacité résiduelle.
\end{proof}

Numériquement,
\[
 \boxed{\lambda_*=0.9892859130191415\ldots},
 \qquad
 0<\lambda_*<1.
\]
La positivité découle déjà de
\(0<\Delta_4<10^{-3}\) et \(5/468<1\). Elle ne dépend pas d’une comparaison
métrologique.

\section{Foncteur d’action positive sur la catégorie des chemins compatibles}

Rappelons le graphe orienté \(\Gtr\) et sa catégorie libre
\(\mathsf P(\Gtr)\). Pour une arête
\(e:x\to y\), nous définissons
\[
 g(e)=\mathbf1_{\{d_y\ne d_x\}},
 \qquad
 s(e)=\mathbf1_{\{m_y\ne m_x\}},
\]
où \(g\) enregistre une rotation et \(s\) un changement de feuillet. La flèche du
scalaire \(\lambda_*\) vers les arêtes ne se déduit pas du
théorème des traces. Elle est déclarée par la clause \HLoc : une rotation a un coût
unitaire, un changement de feuillet a pour coût \(\lambda_*\), et les deux indicateurs
s’additionnent lorsqu’ils apparaissent sur la même arête. Sous \HLoc, l’action locale est
\[
 \boxed{a_*(e)=g(e)+\lambda_*s(e).}
\]
Pour une route \(\gamma=e_1\cdots e_n\),
\[
 \mathcal A_*(\gamma)
 =\sum_{j=1}^na_*(e_j)
 =\#\operatorname{giros}(\gamma)
 +\lambda_*\#\operatorname{cambios\ de\ hoja}(\gamma).
\]
Pour la route vide \(\mathbf1_x\), on fixe
\(\mathcal A_*(\mathbf1_x)=0\). Alors
\[
 \boxed{
 \mathcal A_*(\gamma_1\gamma_2)
 =\mathcal A_*(\gamma_1)+\mathcal A_*(\gamma_2).
 }
\]
Par conséquent,
\[
 \mathcal A_*:\mathsf P(\Gtr)\longrightarrow(\mathbb R_{\ge0},+)
\]
est un foncteur monoïdal, ou cocycle additif de chemins.

Les indicateurs \(g\) et \(s\) sont insensibles au sens de parcours.
L’orientation n’est pas perdue : elle réside dans la droite \(L_{\rm or}\) et dans
l’involution étoile, tandis que le coût positif réside dans les
projecteurs. Confondre ces deux données introduirait des signes négatifs dans
l’action.

Pour \(\beta>0\), l’opérateur de transfert associé est
\[
 (\mathcal L_{\beta,*}f)(y)
 =\sum_{e:x\to y}e^{-\beta a_*(e)}f(x).
\]
Tous ses poids sont strictement positifs. Comme le potentiel dépend seulement
des deux indicateurs locaux, il commute avec l’action grossière
\((\mathbb Z/3\mathbb Z)^3\) et respecte la décomposition en \(27\)
secteurs de Bloch utilisée dans la section suivante.

Une dynamique déterministe finie non pondérée agit par des matrices de
permutation et possède des valeurs singulières égales à un. Le cocycle
\(\mathcal A_*\) transforme l’inventaire des chemins en un opérateur de
transfert qui préserve le cône positif, c’est-à-dire en une matrice à
poids strictement positifs sur les arêtes autorisées. Cette propriété ne
démontre pas à elle seule la contraction, la simplicité de la valeur propre dominante ni
la séparation spectrale ; ces conclusions exigent une normalisation ou un
argument spectral supplémentaire.

\ifdefined\HMTInternalTraceability
\section{Induction par le quotient et extension dynamique}

L’expression \emph{induite par le noyau topologique de phase enrichi} possède deux sens qui doivent rester
séparés.

\begin{enumerate}
  \item \textbf{Induction par l’objet quotient.} L’émetteur exécutable
  détermine \(F\), son facteur polaire \(J_F\), l’espace \(H_U\) et la
  trace normalisée. L’état enrichi apporte
  \(\Pi_\pi^{(468)}\), \(P_3\),
  \(K\), \(B_K\), le repère de Witt et \(\Delta_4\). En ce sens
  catégoriel, \(D_{\rm switch}\) se construit uniquement avec des données internes
  APP--TRIT--TPK\@.
  \item \textbf{Induction par une transition totale.} Cette lecture exigerait
  une transition finie exécutable sur toutes les coordonnées de
  \(\mathcal X_{\mathrm{TPK}}^{\mathrm{enr}}\) et une preuve que sa
  mesure image dynamique
  préserve \(J_FH_U\) et produit le potentiel \(a_*\). Cette transition totale
  n’appartient pas au domaine exécutable construit dans cette section.
\end{enumerate}

Le \cref{p3-20260826:thm:ley-accion-hmt} et la construction des chemins closent la première
lecture au moyen de cinq clauses explicites du lecteur :
\begin{description}
  \item[\HTr --- trace du quotient.] Lorsque la multiplicité de présentation
  reste dans la fibre oubliée, une émission distincte est un atome et
  l’état du catalogue est la trace normalisée de son algèbre finie.
  \item[\HAdd --- cocycle d’action.] Le défaut d’une route est le caractère
  positif, normalisé bifactoriellement, additif et invariant du produit
  régional--dodécaphasique.
  \item[\HDef --- défaut de changement de feuillet.] Le second argument du caractère est
  \(\Delta_4P_3\).
  \item[\HRes --- capacité résiduelle.] La capacité autoduale est normalisée à
  un et le défaut la réduit selon \(\lambda=1-\delta\).
  \item[\HLoc --- action locale.] Sur les arêtes de \Gtr, on fixe
  \(a(e)=g(e)+\lambda s(e)\), et l’action d’une route est la somme sur ses
  arêtes.
\end{description}
\HTr impose \(5/468\). \HAdd impose la somme des traces et exclut les termes
mixtes. \HDef fixe le scalaire qui occupe la coordonnée de commutation. \HRes transforme
le défaut en capacité résiduelle et \HLoc transporte cette capacité vers les arêtes.
Ainsi, \(\delta_*\) est une conséquence de \HTr--\HAdd--\HDef, \(\lambda_*\) l’est de
\HTr--\HAdd--\HDef--\HRes et \(\mathcal A_*\) l’est de
\HTr--\HAdd--\HDef--\HRes--\HLoc. \HTr et \HAdd possèdent des caractérisations universelles ; \HDef, \HRes et
\HLoc sont des identifications constitutives internes ; elles ne se déduisent pas de la
transition finie précédente.

\section{Familles compatibles avec les symétries partielles}

\subsection{Symétrie du cycle \(1/2\) et famille compatible de mesures}

L’échange des deux moitiés d’une émission décompose les \(468\)
émissions en \(234\) paires. La mesure uniforme du catalogue comme les
mesures images des germes sont invariantes sous cet échange. Il en va de même
pour toute attribution constante sur chaque paire. L’invariance du retour,
à elle seule, ne prouve pas l’uniformité.

\subsection{Symétrie de \texorpdfstring{\Gtr}{Gamma tr} et famille de coûts}

Les \(7776\) arêtes de \Gtr se répartissent en \(288\) orbites de taille
\(27\) sous \((\mathbb Z/3\mathbb Z)^3\). Les \(2160\) arêtes de
changement de feuillet forment \(80\) orbites :
\[
 80=5\cdot16.
\]
La symétrie permet d’attribuer des coûts différents à ces orbites. Le coût
scalaire commun provient de \HAdd, et non de la cardinalité.

\subsection{Liberté scalaire de la coordonnée \(2^{\mathrm a}\)}

Pour tout \(c\ge0\), l’opérateur
\[
D_c=\Pi_\pi^{(468)}\otimes I_{11}
+cI_{468}\otimes P_3
\]
est positif et le caractère unique de \HAdd donne
\[
\tau(D_c)=\frac5{468}+\frac3{11}c.
\]
La formule interne de \(\Delta_4\) construit un candidat distingué, mais
la propriété universelle de la trace ne démontre pas à elle seule
\(c=\Delta_4\). Cette égalité est précisément le contenu de \HDef.

\subsection{Positivité, linéarité et termes d’ordre \(2\)}

Les fonctions
\[
 1-u-v+cuv,\qquad(1-u)(1-v),\qquad e^{-(u+v)}
\]
peuvent partager la normalisation et le premier ordre. L’absence de \(uv\) se
déduit de l’additivité exacte. Sans \HAdd, la loi linéaire ne serait pas unique.

\section{Hypothèses d’existence, de positivité, de naturalité et d’extension dynamique}

Le résultat dépend des conditions vérifiables suivantes :
\begin{enumerate}
  \item le catalogue ne contient pas \(468\) émissions ou la microfibre n’a pas
  rang cinq ;
  \item le recensement des fibres n’est pas
  \(432\times144,\ 18\times432,\ 18\times1944\) ;
  \item \(J_F\) n’est pas isométrique ou n’entrelace pas les observables du quotient ;
  \item l’un des \(P_G,P_3,Q_3\) cesse d’être un projecteur de rang trois ;
  \item il existe un autre caractère positif, normalisé bifactoriellement, additif et
  invariant qui n’est pas la somme des traces ;
  \item une transition totale du noyau topologique de phase enrichi, une fois construite, induit un
  coefficient de commutation différent de \(\Delta_4\) ;
  \item \(D_{\rm switch}\) acquiert une valeur propre négative ;
  \item l’action accumulée cesse d’être additive sous concaténation.
\end{enumerate}

Les certificats reconstruisent les fibres, les deux mesures, le facteur polaire,
le spectre du défaut, les orbites d’arêtes et les ablations. La valeur de
\(\lambda_*\) apparaît à la fin comme conséquence de ces objets.
\fi
